\documentclass[12pt]{article}

\usepackage[T1]{fontenc}
\usepackage{lmodern}
\usepackage{amsmath,amssymb,amsthm}
\usepackage{tikz-cd}
\usepackage{tikz}
\usepackage{hyperref}
\usepackage{cleveref}
\usepackage{graphicx}
\usepackage{xcolor}
\usepackage{everypage}
\usepackage[margin=1in]{geometry}
\usepackage{enumitem}
\usepackage{mathtools}

\usetikzlibrary{arrows,decorations.pathmorphing,positioning}

\hypersetup{
  colorlinks=true,
  linkcolor=blue!60!black,
  citecolor=blue!60!black,
  urlcolor=blue!60!black
}

\theoremstyle{plain}
\newtheorem{theorem}{Theorem}[section]
\newtheorem{proposition}[theorem]{Proposition}
\newtheorem{lemma}[theorem]{Lemma}
\newtheorem{corollary}[theorem]{Corollary}

\theoremstyle{definition}
\newtheorem{definition}[theorem]{Definition}
\newtheorem{example}[theorem]{Example}

\theoremstyle{remark}
\newtheorem{remark}[theorem]{Remark}

\newcommand{\Cat}{\mathbf{Cat}}
\newcommand{\Set}{\mathbf{Set}}
\newcommand{\Ham}{\mathbf{Ham}}
\newcommand{\Hilb}{\mathbf{Hilb}}
\newcommand{\FHilb}{\mathbf{FHilb}}
\newcommand{\Rep}{\mathbf{Rep}}
\newcommand{\Vect}{\mathbf{Vect}}
\newcommand{\Phys}{\mathbf{Phys}}
\newcommand{\Info}{\mathbf{Info}}
\newcommand{\Phase}{\mathbf{Phase}}
\newcommand{\BG}{\mathrm{B}G}
\newcommand{\Bord}{\mathbf{Bord}}
\newcommand{\Cob}{\mathbf{Cob}}
\newcommand{\MTC}{\mathbf{MTC}}
\newcommand{\Tr}{\mathrm{Tr}}
\newcommand{\id}{\mathrm{id}}

\title{Law II --- Phase-bound Matter:\\
  Functorial Classification of Thermodynamic and Topological Phases}
\author{MagnetonIO Research}
\date{April 2026}

% --- GrokRxiv-style sidebar (Stage 6) -----------------------------------
\AddEverypageHook{%
  \begin{tikzpicture}[remember picture,overlay]
    \node[anchor=west,xshift=2pt,yshift=0pt,
          rotate=90,
          font=\scriptsize\sffamily\color{gray!70}]
      at (current page.west)
      {GrokRxiv \textbullet\ Emergent Spacetime Dynamics \textbullet\ Law II / IV \textbullet\ MagnetonIO Research \textbullet\ \today};
  \end{tikzpicture}%
}
% -----------------------------------------------------------------------

\begin{document}

\maketitle

\begin{abstract}
We present the second installment of the modular research series \emph{Emergent Spacetime
Dynamics}. Building \emph{strictly modularly} on the categorical primitives established
in Law~I --- the matter--information functor $M:\Phys\to\Info$, the symmetric monoidal
structure of physical composition, and the sheaf semantics of local observables --- we
classify phases of matter as functorial equivalence classes. We first formalise the
Landau paradigm as the connected-component invariant
$\pi_{0}([\BG,\Ham_{0}]^{\mathrm{eq}})$, where $\Ham_{0}$ is the groupoid of gapped local
Hamiltonians under gap-preserving adiabatic equivalence and the superscript ``eq''
restricts to the on-site $G$-action. We then
generalise beyond Landau by showing that topological order in $(2{+}1)$D corresponds to
unitary modular tensor categories (MTCs), arising as Drinfeld centers
$\mathcal{Z}(\mathcal{C})$ of unitary fusion categories. Symmetry-protected topological
(SPT) phases are classified by group cohomology $H^{d+1}(G,U(1))$, identified with the
homotopy classes of pointed maps $[\BG,\mathrm{B}^{d+1}U(1)]$. The topological
entanglement entropy
$\gamma=\log\mathcal{D}$ is constructed as the value of an entanglement-entropy functor
on a connected disk; we show it is invariant under the equivalence relation of finite-depth
local circuits. We provide worked computations for the toric code, the SSH chain, and the
Levin--Wen string-net model, and we expose the precise composition hooks that Law~III
will consume in order to lift these classifications to the periodically driven setting.
This paper is part of an explicitly \emph{modular} framework: each law is a self-standing
mathematical layer that composes functorially onto the prior laws, producing emergent
properties at each level of composition. We do \emph{not} claim a unified theory.
\end{abstract}

\tableofcontents

\section{Introduction}
\label{sec:intro}

\subsection{Position in the modular series}

This paper is Paper~II of four in the modular research series \emph{Emergent Spacetime
Dynamics}. The series is organised as a sequence of categorical layers, each composing
on the previous via an explicit functorial lifting:
\[
\text{Law I} \;\xrightarrow{\,\mathrm{Lift}_{\mathrm{I}\to\mathrm{II}}\,}\;
\text{Law II} \;\xrightarrow{\,\mathrm{Lift}_{\mathrm{II}\to\mathrm{III}}\,}\;
\text{Law III} \;\xrightarrow{\,\mathrm{Lift}_{\mathrm{III}\to\mathrm{IV}}\,}\;
\text{Law IV.}
\]
Law~I (\emph{Mathematical Formalisms}, MagnetonIO Research) established the categorical
grammar that all subsequent laws share: symmetric monoidal categories, dagger-compact
closed categories for quantum systems, sheaves and topoi for local observables, operads
and factorisation algebras for QFT-style composition, the matter--information functor
$M:\Phys\to\Info$, and a type-theoretic encoding of these structures.
The present Paper~II uses these primitives to classify phases of matter --- both
thermodynamic and topological --- as functorial equivalence classes. We work entirely
in the language of Law~I.

We emphasise that the framework is \emph{modular}, not unified. Each law contributes one
self-standing mathematical layer; emergent properties at each layer arise from the
\emph{composition} of layers, not from any single overarching theory.

\subsection{One-page recap of Law I}
\label{sec:lawI-recap}

For convenience we collect the Law~I primitives we will need, fixing notation and
references for the remainder of the paper.

\begin{description}[leftmargin=1.5em,style=nextline]
  \item[(L1.1) Symmetric monoidal category]
  A symmetric monoidal category $(\mathcal{C},\otimes,I,\alpha,\lambda,\rho,\sigma)$
  consists of a category $\mathcal{C}$, a tensor functor
  ${\otimes:\mathcal{C}\times\mathcal{C}\to\mathcal{C}}$, a unit object $I$, and natural
  isomorphisms (associator $\alpha$, unitors $\lambda,\rho$, symmetry $\sigma$)
  satisfying Mac~Lane's pentagon and triangle axioms together with the symmetry
  hexagons. Mac~Lane's coherence theorem guarantees that all rebracketings commute, so
  diagrams of structural isomorphisms can be elided in calculations.

  \item[(L1.2) Dagger-compact closed structure]
  A $\dagger$-category is a category equipped with an involutive contravariant identity-
  on-objects functor $\dagger:\mathcal{C}^{\mathrm{op}}\to\mathcal{C}$. A
  $\dagger$-compact closed category additionally has duals
  $A\mapsto A^{*}$ with unit $\eta_{A}:I\to A^{*}\otimes A$ and counit
  $\varepsilon_{A}:A\otimes A^{*}\to I$. The category $\FHilb$ of finite-dimensional
  Hilbert spaces with linear maps is the canonical example. Traces, partial traces and
  the Born rule arise from $(\eta,\varepsilon)$.

  \item[(L1.3) Matter--information functor]
  Law~I posits a symmetric monoidal functor
  $M:\Phys\to\Info$
  from a category of physical systems to a category of information-theoretic structures
  (sets of states, density matrices, classical channels, quantum channels). Whenever a
  physical category has a composition (tensor product), $M$ preserves it.
  Concretely, the prototypical instance of $M$ sends a quantum system with Hilbert
  space $\mathcal{H}\in\FHilb$ to its space of density matrices
  $M(\mathcal{H})=\{\rho\in\mathrm{End}(\mathcal{H}) : \rho\geq 0,\;\Tr\rho=1\}$,
  regarded as the object of statistical states; on morphisms (unitary or completely
  positive maps) $M$ acts by push-forward. We use $M$ in the present paper to translate
  ``Hilbert space of microstates'' (an object of $\Phys$) into ``probability simplex of
  accessible macrostates'' (an object of $\Info$).

  \item[(L1.4) Sheaf semantics of local observables]
  Law~I formalises local observables as a sheaf
  $\mathcal{O}:\mathrm{Open}(X)^{\mathrm{op}}\to\mathbf{Alg}$ assigning algebras of
  operators to open regions. We use this to formulate locality conditions for order
  parameters and entanglement-entropy regions in this paper.

  \item[(L1.5) Functorial semantics]
  Following Lawvere, Law~I treats algebraic theories as monoidal categories whose
  product-preserving functors into a base category are models. We instantiate this in
  \cref{sec:landau} by viewing a Landau theory as a functor $\BG\to\Ham$.
\end{description}

When we cite ``(L1.k)'' below we are using the Law~I primitive of the same number.

\subsection{Outline}

The paper is organised as follows. In \cref{sec:phases-as-classes} we define the category
$\Ham$ of gapped Hamiltonians on a fixed lattice and introduce the equivalence relation
of \emph{gap-preserving adiabatic deformation}; phases are equivalence classes (i.e.\
connected components of $\Ham$). We then define the functor pullback that classifies
phases under symmetry. In \cref{sec:landau} we recover Landau's order-parameter
paradigm in this language. In \cref{sec:spt} we treat SPT phases via group cohomology and
identify the classifying functor with a generalised cohomology functor.
\Cref{sec:mtc} sets up unitary modular tensor categories as the data of a $(2{+}1)$D
topological phase, with Kitaev's quantum double and the Levin--Wen construction
as worked examples. \Cref{sec:eee} defines an entanglement-entropy functor and
proves the Kitaev--Preskill formula for $\gamma=\log\mathcal{D}$ in our setting.
\Cref{sec:examples} works through three canonical examples: the toric code, the SSH
chain, and the Fibonacci Levin--Wen model. \Cref{sec:law3-hooks} states the precise
composition hooks Law~III will consume. \Cref{sec:open} lists open problems, and
\cref{sec:conclusion} concludes.

\section{Mathematical Framework: Phases as Equivalence Classes under Functor Pullback}
\label{sec:phases-as-classes}

\subsection{The category of gapped Hamiltonians}

Fix a (regular, finite or countable) lattice $\Lambda$ and a finite-dimensional on-site
Hilbert space $\mathfrak{h}$. Following Law~I (L1.2), the total Hilbert space is the
infinite tensor product (or a thermodynamic limit thereof)
$\mathcal{H}_{\Lambda}=\bigotimes_{x\in\Lambda}\mathfrak{h}$, an object of (a suitable
completion of) $\FHilb$.

\begin{definition}[Local Hamiltonian]
\label{def:local-ham}
A \emph{local Hamiltonian} on $\Lambda$ is a sum
\[
H \;=\; \sum_{X\subset\Lambda} h_{X}, \qquad h_{X}\in\mathrm{End}(\mathcal{H}_{X}),
\]
where each $h_{X}$ is supported on a finite region $X\subset\Lambda$ with diameter
bounded by a fixed range $r$, and the operator norm of $h_{X}$ is uniformly bounded.
\end{definition}

\begin{definition}[Gap]
\label{def:gap}
A local Hamiltonian $H$ is \emph{gapped} with gap $\Delta>0$ if its spectrum on
$\mathcal{H}_{\Lambda}$ has a single (possibly degenerate) ground subspace separated
from the rest of the spectrum by at least $\Delta$, uniformly as $|\Lambda|\to\infty$.
\end{definition}

\begin{definition}[Category $\Ham$]
\label{def:ham-cat}
The category $\Ham=\Ham(\Lambda,\mathfrak{h})$ has:
\begin{itemize}
  \item objects: gapped local Hamiltonians on $\Lambda$;
  \item morphisms $H_{0}\to H_{1}$: gap-preserving piecewise-smooth paths
    $\{H_{s}\}_{s\in[0,1]}$ of local Hamiltonians with $H_{s=0}=H_{0}$, $H_{s=1}=H_{1}$,
    and $\inf_{s}\Delta(H_{s})>0$;
  \item composition: concatenation of paths;
  \item identity: the constant path.
\end{itemize}
\end{definition}

By construction $\Ham$ is a groupoid (every path is reversible), and morphisms are
homotopic iff they are connected through a homotopy of gap-preserving paths. We will
write $\pi_{0}(\Ham)$ for the set of connected components.

\begin{remark}
$\Ham$ is symmetric monoidal: the tensor product
$H \boxtimes H'$ of two Hamiltonians on disjoint lattices is again a gapped local
Hamiltonian on the union lattice; gap-preserving paths tensor componentwise. This is the
direct image under the matter--information functor $M$ (L1.3) of the symmetric monoidal
structure on $\Phys$.
\end{remark}

\subsection{Phases as functorial equivalence classes}

\begin{definition}[Phase]
\label{def:phase}
A \emph{phase} of matter (in the present setting) is a connected component of $\Ham$:
\[
\Phase \;=\; \pi_{0}(\Ham).
\]
Two Hamiltonians belong to the same phase iff they are connected by a gap-preserving
adiabatic path.
\end{definition}

\begin{remark}
Three points deserve emphasis. (i) $\Phase$ is a set, not a category --- the phase
classification forgets the choice of path. (ii) The thermodynamic limit
$|\Lambda|\to\infty$ is essential: at finite volume every two Hamiltonians can be
connected by some path, but the gap typically closes uniformly only on a measure-zero
subset. (iii) The lattice $\Lambda$ and on-site Hilbert space $\mathfrak{h}$ are part
of the data; phases are defined relative to a microscopic site structure.
\end{remark}

\begin{definition}[Symmetry action]
\label{def:sym-action}
Let $G$ be a group acting on $\mathfrak{h}$ by a unitary representation
$\rho:G\to U(\mathfrak{h})$. The induced \emph{on-site} action on $\mathcal{H}_{\Lambda}$
is $U_{g}=\bigotimes_{x\in\Lambda}\rho(g)$. We say $H\in\Ham$ is \emph{$G$-symmetric}
if $[U_{g},H]=0$ for all $g\in G$.
\end{definition}

\begin{definition}[Category $\Ham_{G}$]
\label{def:HamG}
The category $\Ham_{G}$ has $G$-symmetric Hamiltonians as objects, and gap-preserving
paths $\{H_{s}\}$ \emph{through $G$-symmetric Hamiltonians} as morphisms.
\end{definition}

The forgetful functor $U:\Ham_{G}\to\Ham$ remembers only the underlying Hamiltonian
without its $G$-equivariance. The induced map on connected components
$U_{*}:\pi_{0}(\Ham_{G})\to\pi_{0}(\Ham)$ is in general neither injective nor
surjective: phases that are distinct under $G$-symmetry can become equivalent when the
symmetry is forgotten (this is precisely the SPT phenomenon, \cref{sec:spt}).

\subsection{The classifying space of $G$ as a one-object category}

For a group $G$, write $\BG$ for the one-object groupoid with morphism set $G$ and
composition given by group multiplication. A functor $\BG\to\mathcal{D}$ for any category
$\mathcal{D}$ picks out an object $D$ and a homomorphism $G\to\mathrm{Aut}_{\mathcal{D}}(D)$.

\begin{proposition}[Symmetric Hamiltonians as $\BG$-functors]
\label{prop:hamG-as-functor}
Let $\Ham_{0}$ denote the homotopy 1-category of $\Ham$: its objects are gapped local
Hamiltonians on $(\Lambda,\mathfrak{h})$, and its morphisms are \emph{homotopy classes}
of gap-preserving paths. Explicitly,
$\mathrm{Hom}_{\Ham_{0}}(H,H')$ is the set of path-homotopy classes
$[\{H_{s}\}_{s\in[0,1]}]$ where $H_{s=0}=H$, $H_{s=1}=H'$, and $\inf_{s}\Delta(H_{s})>0$,
under the equivalence relation that two paths are identified iff they are homotopic
through gap-preserving paths fixing endpoints. Composition is concatenation of paths.
With this definition $\mathrm{Hom}_{\Ham_{0}}(H,H)$ is a group (the fundamental group of
the gapped locus at the basepoint $H$). Note that $\Ham_{0}=\pi_{\leq 1}(\Ham)$ is the
$1$-truncation of the full $\infty$-groupoid $\Ham$; we will not need higher homotopy
in this paper. Let
$[\BG,\Ham_{0}]^{\mathrm{eq}}$ denote the full subcategory of the functor category
$[\BG,\Ham_{0}]$ consisting of those functors $F$ for which the induced homomorphism
$G\to\mathrm{Aut}_{\Ham_{0}}(F(\ast))$ factors through the homomorphism
$\rho_{\mathrm{ad}}:G\to\mathrm{Aut}_{\Ham_{0}}(F(\ast))$ given by adjoint action
$\rho_{\mathrm{ad}}(g)=\mathrm{Ad}_{U_{g}}$ where $\mathrm{Ad}_{U_{g}}(H)=U_{g}H
U_{g}^{\dagger}$. (For $G$-symmetric Hamiltonians this adjoint action is the identity
loop; the data of the functor is precisely a recording of which on-site action $U_{g}$
is being used.) Then
\[
\Ham_{G} \;\simeq\; [\BG,\Ham_{0}]^{\mathrm{eq}},
\]
as groupoids.
\end{proposition}

\begin{proof}
We construct mutually inverse functors and check that they are equivalences.

\emph{Functor $\Phi:\Ham_{G}\to[\BG,\Ham_{0}]^{\mathrm{eq}}$.}
Given a $G$-symmetric Hamiltonian $H$ (so $[U_{g},H]=0$ for all $g\in G$), define a
functor $F_{H}:\BG\to\Ham_{0}$ by setting $F_{H}(\ast)=H$ and, for each $g\in G$,
$F_{H}(g):=[\mathrm{Ad}_{U_{g}}]\in\mathrm{Aut}_{\Ham_{0}}(H)$, the homotopy class of
the loop $s\mapsto U_{g(s)}HU_{g(s)}^{\dagger}$ where $g(s)$ is any path in
$U(\mathfrak{h})$ from $\mathbb{1}$ to the on-site representative $U_{g}$. Two
observations make this well-defined and functorial.
\emph{(a)} Because $[U_{g},H]=0$ we have $U_{g}HU_{g}^{\dagger}=H$ on the nose, so the
endpoint of the loop coincides with the basepoint $H$. \emph{(b)} The on-site action
gives a group homomorphism $G\to U(\mathfrak{h})^{\otimes\Lambda}$, so
$\mathrm{Ad}_{U_{gh}}=\mathrm{Ad}_{U_{g}}\circ\mathrm{Ad}_{U_{h}}$ as automorphisms of
$\Ham_{0}$; hence the assignment $g\mapsto[\mathrm{Ad}_{U_{g}}]$ is a homomorphism
$G\to\mathrm{Aut}_{\Ham_{0}}(H)$, factoring through $\rho_{\mathrm{ad}}$ as required.
This is the precise meaning of ``factoring through $\rho_{\mathrm{ad}}$.''
On morphisms (gap-preserving $G$-symmetric paths $\{H_{s}\}$), $\Phi$ sends $\{H_{s}\}$
to the natural transformation $\Phi(\{H_{s}\}):F_{H_{0}}\Rightarrow F_{H_{1}}$ whose
single component at $\ast$ is the homotopy class of the path itself. Naturality means
the square
\[
\begin{tikzcd}[column sep=large]
H_{0} \arrow[r,"{[\mathrm{Ad}_{U_{g}}]}"] \arrow[d,"\{H_{s}\}"'] & H_{0} \arrow[d,"\{H_{s}\}"] \\
H_{1} \arrow[r,"{[\mathrm{Ad}_{U_{g}}]}"] & H_{1}
\end{tikzcd}
\]
commutes in $\Ham_{0}$ for every $g\in G$. This is exactly the condition that
$\{H_{s}\}$ is a $G$-symmetric path, since both compositions equal the path
$\{H_{s}\}$ itself --- the top and bottom horizontal arrows are constant loops at the
respective endpoints because each $H_{s}$ commutes with $U_{g}$.

\emph{Functor $\Psi:[\BG,\Ham_{0}]^{\mathrm{eq}}\to\Ham_{G}$.}
Given $F\in[\BG,\Ham_{0}]^{\mathrm{eq}}$, set $\Psi(F)=F(\ast)$. By hypothesis $F(g)=
[\mathrm{Ad}_{U_{g}}]$ in $\mathrm{Aut}_{\Ham_{0}}(F(\ast))$, i.e.\ the homotopy class
of the adjoint loop $H_{s}=U_{g(s)}F(\ast)U_{g(s)}^{\dagger}$ for some choice of
representative path. Since this loop is the constant loop in the homotopy category
$\Ham_{0}$, we have $\mathrm{Ad}_{U_{g}}(F(\ast))=U_{g}F(\ast)U_{g}^{\dagger}$
gap-equivalent to $F(\ast)$ for all $g\in G$. Combined with strict equality at the
identity element $g=e$, the cocycle condition gives
$U_{g}F(\ast)U_{g}^{\dagger}=F(\ast)$, i.e.\ $[U_{g},F(\ast)]=0$, so
$F(\ast)\in\Ham_{G}$. On a natural transformation $F\Rightarrow F'$, $\Psi$ takes
its component at $\ast$.

\emph{Inverse equivalence.}
$\Psi\circ\Phi=\id_{\Ham_{G}}$ on the nose. For $\Phi\circ\Psi$, given $F$, the
composite functor $F_{F(\ast)}$ has the same object value $F(\ast)$ and the same
morphism values up to the on-site identification; the natural isomorphism
$F\cong F_{F(\ast)}$ has identity component at $\ast$, so the two functors are equal up
to canonical isomorphism. This yields a natural equivalence $\Phi\circ\Psi\cong\id$.
\end{proof}

\begin{remark}
The condition that the homomorphism $G\to\mathrm{Aut}_{\Ham_{0}}(F(\ast))$ factor
through the on-site subgroup is essential. Concretely, the \emph{on-site subgroup} of
$\mathrm{Aut}_{\Ham_{0}}(H)$ is the image of $\rho_{\mathrm{ad}}:G\to
\mathrm{Aut}_{\Ham_{0}}(H)$, which by definition consists of those automorphisms of $H$
realised by adjoint action of unitaries of the form $\bigotimes_{x\in\Lambda}\rho(g)$
(with $\rho:G\to U(\mathfrak{h})$ a fixed on-site representation). An arbitrary functor
$\BG\to\Ham_{0}$ would encode an arbitrary --- not necessarily on-site --- $G$-action
on the Hamiltonian, and such functors classify projective representations and
anomalies, which we treat in \cref{sec:spt} via group cohomology.
\end{remark}

\begin{remark}
\Cref{prop:hamG-as-functor} is the formal mechanism by which Law~I's functorial
semantics (L1.5) classifies phases. The Landau paradigm is the special case where the
target is the category of Hamiltonians built from a fixed local order-parameter degree
of freedom; topological phases correspond to functors landing on Hamiltonians whose
ground states are not products in any local basis.
\end{remark}

\subsection{Phases under symmetry as functor pullbacks}

\begin{definition}[Symmetric phase]
\label{def:sym-phase}
A \emph{$G$-symmetric phase} is an element of $\pi_{0}(\Ham_{G})$.
\end{definition}

The map $U_{*}:\pi_{0}(\Ham_{G})\to\pi_{0}(\Ham)$ has a kernel and cokernel that
encode, respectively, SPT phases (kernel) and spontaneous symmetry breaking (cokernel).
This is one of the two main classification tools we use; the other is the assignment of
algebraic invariants to ground states (modular data of an MTC, group-cohomology classes,
etc.) which we develop in subsequent sections.

\begin{theorem}[Phase-classification exact sequence]
\label{thm:phase-classification}
Let $\pi_{0}(\Ham_{G})^{\mathrm{SRE}}\subseteq\pi_{0}(\Ham_{G})$ denote the
\emph{short-range entangled} $G$-symmetric phases, i.e.\ those phases whose ground
state can be transformed into a product state by a finite-depth local quantum circuit
that need not preserve $G$. Let $\pi_{0}(\Ham)^{\mathrm{SRE}}$ be the corresponding set
of (non-symmetric) SRE phases; this set has a single element, the trivial phase
$\ast$, by Chen--Gu--Wen \cite{ChenGuLiuWen2013}. The forgetful map
$U_{*}:\pi_{0}(\Ham_{G})\to\pi_{0}(\Ham)$ restricts to
$U_{*}^{\mathrm{SRE}}:\pi_{0}(\Ham_{G})^{\mathrm{SRE}}\to\pi_{0}(\Ham)^{\mathrm{SRE}}=\{\ast\}$,
and the fibre $(U_{*}^{\mathrm{SRE}})^{-1}(\ast)$ is exactly the set of SPT phases
$\mathrm{SPT}^{d}(G)$ defined in \cref{def:spt}. Concretely we have a sequence of
pointed sets
\[
\mathrm{SPT}^{d}(G) \;\xhookrightarrow{}\; \pi_{0}(\Ham_{G})^{\mathrm{SRE}}
\;\xrightarrow{\,U_{*}^{\mathrm{SRE}}\,}\; \pi_{0}(\Ham)^{\mathrm{SRE}} = \{\ast\},
\]
which is exact in the sense that the kernel of $U_{*}^{\mathrm{SRE}}$ (the preimage of
the basepoint $\ast$) is exactly the image of the inclusion. Adding spontaneously
symmetry-broken phases extends the sequence: the cokernel of the inclusion
$\mathrm{SPT}^{d}(G)\hookrightarrow\pi_{0}(\Ham_{G})^{\mathrm{SRE}}$ classifies SSB
patterns by subgroups $H\subseteq G$.
\end{theorem}

\begin{proof}
Restriction to SRE phases is well-defined because the property of being SRE is preserved
by gap-preserving deformations \cite{ChenGuLiuWen2013}. By definition of SRE, every SRE
phase becomes the trivial phase upon forgetting symmetry, so
$U_{*}^{\mathrm{SRE}}$ lands in $\{\ast\}$.

The inclusion $\mathrm{SPT}^{d}(G)\hookrightarrow\pi_{0}(\Ham_{G})^{\mathrm{SRE}}$ is by
definition (\cref{def:spt}). Surjectivity onto the kernel: any
$[H]\in\pi_{0}(\Ham_{G})^{\mathrm{SRE}}$ with $U_{*}[H]=\ast$ has, upon forgetting
symmetry, a representative connected to a product state by a (non-symmetric) finite-depth
circuit. By definition this is an SPT phase. Conversely, any SPT phase has SRE ground
state (no anyons) and is in the kernel of $U_{*}$ by construction.

For SSB patterns, take $[H]\in\pi_{0}(\Ham_{G})$ whose ground state is a $G$-symmetric
superposition of states each individually invariant only under a subgroup $H\subseteq G$.
Modding out by SRE phases identifies $[H]$ with the data of the unbroken subgroup $H$;
this is the standard Goldstone classification \cite{LandauLifshitz1980}.
\end{proof}

\begin{remark}
This theorem is best regarded as exhibiting an exact sequence of pointed sets in the
homotopy category of phase classifications; we have not asserted higher homotopical
content. A topos-theoretic refinement, in which one carries the full groupoid structure
of $\Ham_{G}$ rather than just $\pi_{0}$, gives an analogous fibre sequence at the
level of $\infty$-groupoids and is the natural setting for higher-dimensional
generalisations. We do not pursue this here.
\end{remark}

\section{Landau Theory in Categorical Form}
\label{sec:landau}

Before passing to topological phases we recover the Landau paradigm in the categorical
language of \cref{sec:phases-as-classes}. This serves three purposes: it shows the
present framework is at least as expressive as classical phase theory; it fixes notation
that the topological theory generalises; and it isolates the precise feature of Landau
theory --- representability of an order parameter by a local section --- that
topological order violates.

\subsection{Order parameters as local sections}

\begin{definition}[Order parameter representation]
\label{def:op-rep}
Let $G$ act on $\mathfrak{h}$. An \emph{order-parameter representation} is a
finite-dimensional unitary representation $V$ of $G$ together with a $G$-equivariant
local map $\Phi_{x}:\mathrm{End}(\mathfrak{h})\to V$ for each site $x\in\Lambda$, all
related by translation.
\end{definition}

In a state $\omega$ on $\mathcal{H}_{\Lambda}$, the order parameter is
$\langle\Phi_{x}\rangle_{\omega}=\omega(\Phi_{x})\in V$.

\begin{definition}[Landau functor]
\label{def:landau-functor}
The \emph{Landau functor} of a $G$-symmetric Hamiltonian $H$ is
\[
\mathcal{L}_{H}:\BG\to\Vect, \qquad
\mathcal{L}_{H}(\ast) = V, \qquad
\mathcal{L}_{H}(g) = \rho_{V}(g),
\]
where $V$ is the order-parameter representation and the assigned map records the action
of $G$ on the expectation $\langle\Phi_{x}\rangle$ in the ground state.
\end{definition}

\begin{proposition}[Landau classification, functorial form]
\label{prop:landau}
Two $G$-symmetric Hamiltonians $H_{0},H_{1}$ are in the same Landau phase iff their
Landau functors $\mathcal{L}_{H_{0}}$ and $\mathcal{L}_{H_{1}}$ are naturally isomorphic
as functors $\BG\to\Vect$. Equivalently, two Hamiltonians are in the same Landau phase
iff their order-parameter expectations lie on the same $G$-orbit (up to $V$-isomorphism).
\end{proposition}

\begin{proof}
A natural isomorphism $\eta:\mathcal{L}_{H_{0}}\Rightarrow\mathcal{L}_{H_{1}}$ is a
$G$-intertwiner between the order-parameter representations. Existence of such an
intertwiner is exactly the condition that the order parameters lie on the same $G$-orbit,
which by Goldstone's theorem and standard Landau analysis is equivalent to lying in the
same gap-preserving phase \cite{LandauLifshitz1980,Wilson1983}.
\end{proof}

\begin{remark}
The Ginzburg--Landau free energy
$
f[\varphi]=a(T)|\varphi|^{2}+b|\varphi|^{4}+c|\nabla\varphi|^{2}+\cdots
$
is recovered as the action of $\mathcal{L}_{H}$ on a generating loop in $\BG$, weighted
by the relevant invariant polynomials in $V$.
\end{remark}

\subsection{Landau is not enough}

The Landau classification is precisely the data
$(V,\langle\Phi\rangle\text{ up to orbit})$. Two limitations are immediate.
\begin{enumerate}[label=(\roman*)]
  \item It is local: it sees only the value of an on-site (or short-range) functional.
  \item It is symmetry-breaking-centric: it distinguishes phases via the residual
    subgroup $H\subseteq G$ stabilising $\langle\Phi\rangle$.
\end{enumerate}
A topological phase is, by definition, one where $\langle\Phi\rangle\equiv 0$ for every
local $\Phi$ and yet the ground subspace has structure (degeneracy, anyonic excitations,
non-trivial entanglement) detectable only nonlocally. The Landau functor of a
topological phase is trivial; we need a richer invariant.

\section{Symmetry-Protected Topological Phases via Group Cohomology}
\label{sec:spt}

\subsection{SPT phases as the SSB-trivial subset}

\begin{definition}[SPT phase]
\label{def:spt}
A \emph{symmetry-protected topological (SPT) phase} of $G$-symmetric Hamiltonians is an
element of $\ker U_{*}\subseteq\pi_{0}(\Ham_{G})$: a $G$-symmetric phase that becomes
trivial upon forgetting $G$-symmetry, but is nontrivial as a $G$-symmetric phase.
\end{definition}

\subsection{Group cohomology}

For a group $G$ and a $G$-module $M$, the standard bar complex computes group
cohomology. The cochain in degree $n$ is the abelian group
$C^{n}(G,M)=\{f:G^{n}\to M\}$, with differential
\begin{align*}
(\delta f)(g_{1},\ldots,g_{n+1}) &= g_{1}\!\cdot f(g_{2},\ldots,g_{n+1}) \\
  &\quad - \sum_{i=1}^{n}(-1)^{i-1} f(g_{1},\ldots,g_{i}g_{i+1},\ldots,g_{n+1}) \\
  &\quad +(-1)^{n}f(g_{1},\ldots,g_{n}).
\end{align*}
We write
$Z^{n}=\ker\delta$, $B^{n}=\mathrm{im}\,\delta$, and
$H^{n}(G,M)=Z^{n}/B^{n}$. We work with $M=U(1)$ throughout, with the trivial action.

\subsection{The Chen--Gu--Liu--Wen classification}

\begin{theorem}[Chen--Gu--Liu--Wen \cite{ChenGuLiuWen2013}]
\label{thm:CGLW}
Let $G$ be a finite group acting on-site, and let $d\geq 1$. The set of bosonic SPT
phases on $\Lambda=\mathbb{Z}^{d}$ with on-site $G$-symmetry is in bijection with
$H^{d+1}(G,U(1))$.
\end{theorem}

The bijection is constructive: given a cocycle $\omega\in Z^{d+1}(G,U(1))$, one builds a
fixed-point Hamiltonian whose ground state is a sum of group elements weighted by
$\omega$ on simplices of a triangulation. We will not reproduce the full construction
here; see \cite{ChenGuLiuWen2013}. The key categorical content for us is:

\begin{proposition}[SPT classification as a functor]
\label{prop:spt-functor}
There is a functor
\[
\mathrm{SPT}^{d}: \mathbf{Grp} \to \mathbf{Ab}, \qquad
G \mapsto H^{d+1}(G,U(1)),
\]
and the assignment $H\mapsto[\omega_{H}]$ refines the phase classification map
\[
\pi_{0}(\Ham_{G}) \;\to\; H^{d+1}(G,U(1))
\]
on the SPT subset.
\end{proposition}

\begin{proof}
$H^{d+1}(-,U(1))$ is functorial in $G$ via pullback along group homomorphisms; the map
on phases is induced by restriction of the symmetry action.
\end{proof}

\begin{example}[Haldane phase]
\label{ex:haldane}
For $G=SO(3)$ and $d=1$ (1D spin chain), one has
\[
H^{2}(SO(3),U(1))\cong\mathbb{Z}_{2}
\]
\cite{ChenGuLiuWen2013}. Here $SO(3)$ is the rotational symmetry of an
isotropic Heisenberg spin-1 chain, e.g.\ the AKLT chain
\begin{equation}
\label{eq:aklt}
H_{\mathrm{AKLT}}=\sum_{i}\Bigl[\,\vec{S}_{i}\!\cdot\vec{S}_{i+1}
+\tfrac{1}{3}(\vec{S}_{i}\!\cdot\vec{S}_{i+1})^{2}\,\Bigr].
\end{equation}
Each site carries the spin-1 representation, and on-site $SO(3)$ acts diagonally on the
chain. The non-trivial element of $H^{2}(SO(3),U(1))$ corresponds to the AKLT/Haldane
phase. Concretely, the relevant $2$-cocycle is the projective representation class of
$SO(3)$ on the boundary spin-$\tfrac{1}{2}$ degree of freedom that emerges at each end
of an open AKLT chain.
\end{example}

\begin{example}[1D fermionic SPT]
\label{ex:1d-fermionic}
The Kitaev chain (1D class D superconductor) is classified by
$H^{2}(\mathbb{Z}_{2}^{f},U(1))\cong\mathbb{Z}_{2}$ in the bosonic-shadow framework, but
the full fermionic classification uses spin cobordism rather than group cohomology
\cite{FreedHopkins2021}. We mention this only to flag that the classification
$H^{d+1}(G,U(1))$ is the bosonic answer; fermionic systems require an upgrade to
generalised cohomology, which is the open-problem item in \cref{sec:open}.
\end{example}

\subsection{The classifying-space picture}

The cohomology group $H^{n}(G,U(1))$ is the set of homotopy classes of pointed maps
\[
H^{n}(G,U(1)) \;\cong\; [\BG, \mathrm{B}^{n} U(1)],
\]
where $\mathrm{B}^{n}U(1)=K(U(1),n)$ is the Eilenberg--MacLane space. We will use this
identification in \cref{sec:law3-hooks} to formulate the lifting to Floquet phases:
Law~III replaces $\BG$ with $\BG\times BS^{1}$ (adding the discrete time-translation
group), and the relevant cohomology becomes $H^{*}(\BG\times BS^{1},U(1))$, which by
K\"unneth decomposes in a manner that exposes Floquet-specific phases.

\section{Modular Tensor Categories and Topological Order}
\label{sec:mtc}

We now leave the Landau and SPT regimes and treat the most informationally rich phases:
those with topological order.

\subsection{Unitary fusion categories}

\begin{definition}[Unitary fusion category (UFC)]
\label{def:ufc}
A UFC $\mathcal{C}$ is a $\mathbb{C}$-linear semisimple monoidal $\dagger$-category
with finitely many isomorphism classes of simple objects $\{X_{i}\}_{i\in I}$, a simple
unit object $\mathbf{1}=X_{0}$, dual objects $X_{i}^{*}$, and unitary
associators (F-matrices) satisfying the Mac~Lane pentagon and unitarity.
\end{definition}

We write $N_{ij}^{k}\in\mathbb{Z}_{\geq 0}$ for the fusion multiplicity:
$X_{i}\otimes X_{j}\cong\bigoplus_{k}N_{ij}^{k}X_{k}$.

\subsection{Modular tensor categories}

\begin{definition}[MTC]
\label{def:mtc}
A \emph{modular tensor category} is a UFC $\mathcal{C}$ equipped with a unitary braiding
$\beta_{X,Y}:X\otimes Y\to Y\otimes X$ and a ribbon twist $\theta_{X}:X\to X$ such that
the modular $S$-matrix
\[
S_{ij} = \frac{1}{\mathcal{D}}\Tr(\beta_{X_{j},X_{i}}\circ\beta_{X_{i},X_{j}}),
\qquad \mathcal{D}=\sqrt{\sum_{i}d_{i}^{2}},
\]
is non-degenerate, where $d_{i}$ is the quantum dimension of $X_{i}$.
\end{definition}

\begin{remark}
The non-degeneracy of $S$ is the modular condition. By a theorem of M\"uger
\cite{Muger2003}, $S$ is non-degenerate iff the M\"uger center of $\mathcal{C}$ is trivial.
\end{remark}

\subsection{Topological order as MTC data}

\begin{theorem}[Wen, Levin--Wen, Kitaev: topological order is MTC data]
\label{thm:top-order-mtc}
Every $(2{+}1)$D topologically ordered gapped phase admits, as ground-state algebraic
data, a unitary modular tensor category $\mathcal{C}$ such that:
\begin{enumerate}[label=(\arabic*)]
  \item Simple objects of $\mathcal{C}$ are anyon types.
  \item Fusion rules are $X_{i}\otimes X_{j}\cong\bigoplus_{k}N_{ij}^{k}X_{k}$.
  \item Braiding statistics are encoded in $\beta$ and $\theta$.
  \item Ground-state degeneracy on a genus-$g$ surface is
    $\mathrm{GSD}(g) = \sum_{i\in I} d_{i}^{2(g-1)}$ (Verlinde formula
    \cite{Wang2010}). At $g=1$ this gives $\mathrm{GSD}(1)=|I|$ (number of simple
    objects, i.e.\ anyon types), valid for abelian \emph{and} non-abelian theories. At
    higher genus the formula involves the quantum dimensions $d_{i}$ explicitly.
\end{enumerate}
Conversely, the Levin--Wen string-net construction \cite{LevinWen2005} produces, from a
unitary fusion category $\mathcal{C}_{0}$, a $(2{+}1)$D commuting-projector lattice
Hamiltonian whose anyon MTC is the Drinfeld center $\mathcal{Z}(\mathcal{C}_{0})$.
\end{theorem}

\subsection{Kitaev's quantum double and the toric code}

For a finite group $G$, the quantum double Hopf algebra is
$D(G)=\mathrm{Fun}(G)\rtimes\mathbb{C}[G]$. Its representation category
$\Rep(D(G))$ is a UMTC.

\begin{example}[Toric code, on-site unitary $\mathbb{Z}_{2}$]
\label{ex:toric}
Here $G=\mathbb{Z}_{2}$ is the on-site \emph{unitary} $\mathbb{Z}_{2}$ symmetry
(``$\mathbb{Z}_{2}$ gauge'' of the underlying lattice qubit; not a time-reversal or
charge symmetry). The category $\Rep(D(\mathbb{Z}_{2}))$ has four simple objects:
$\mathbf{1}$, $e$, $m$, $\varepsilon=em$,
with fusion rules $e\otimes e=m\otimes m=\varepsilon\otimes\varepsilon=\mathbf{1}$,
$e\otimes m=\varepsilon$, etc. All quantum dimensions are $1$, so $\mathcal{D}=2$ and
$\gamma=\log 2$. The braiding of $e$ around $m$ is $-1$ (mutual semions), and
$\varepsilon$ has topological spin $\theta_{\varepsilon}=-1$ (a fermion).
\end{example}

The Kitaev Hamiltonian on the square lattice (qubit on each edge):
\begin{equation}
\label{eq:toric-ham}
H_{\mathrm{TC}} = -\sum_{v}A_{v} - \sum_{p}B_{p}, \qquad
A_{v}=\prod_{e\ni v}\sigma^{x}_{e}, \qquad
B_{p}=\prod_{e\in\partial p}\sigma^{z}_{e},
\end{equation}
realises a representative of the toric-code phase. All terms commute, and the ground
space on a torus is four-dimensional, encoding the $[[n,2,n^{1/2}]]$ stabiliser code that
plays a central role in fault-tolerant quantum computation \cite{Kitaev2003}.

\subsection{Levin--Wen string-nets and Drinfeld centers}

\begin{theorem}[Levin--Wen \cite{LevinWen2005}]
\label{thm:lw}
Given a UFC $\mathcal{C}_{0}$ with simple objects $\{a_{i}\}$, $F$-symbols $F^{abc}_{def}$,
and quantum dimensions $\{d_{i}\}$, define a lattice Hamiltonian on the honeycomb whose
edges carry labels $a_{i}$ and whose plaquette terms project onto the string-net
condensate. The resulting model is exactly solvable, has gapped ground states, and the
modular tensor category of its anyonic excitations is the Drinfeld center
$\mathcal{Z}(\mathcal{C}_{0})$.
\end{theorem}

\begin{example}[Doubled Fibonacci]
\label{ex:fib}
Let $\mathcal{C}_{0}=\mathrm{Fib}$ be the Fibonacci fusion category (two simple objects
$\mathbf{1},\tau$ with $\tau\otimes\tau=\mathbf{1}\oplus\tau$, golden-ratio quantum
dimension $d_{\tau}=\phi$). Then $\mathcal{Z}(\mathrm{Fib})\cong\mathrm{Fib}\boxtimes
\overline{\mathrm{Fib}}$. This is a non-abelian topological order capable of universal
topological quantum computation \cite{Wang2010}.
\end{example}

\subsection{Compositional connection to Law I}

The MTC structure is exactly the Law~I primitive of a braided monoidal category with
duals, decorated by a non-degenerate $S$-matrix. The Drinfeld center
$\mathcal{Z}:\mathbf{FusCat}\to\mathbf{MTC}$ is a functor on the Law~I 2-category of
fusion categories. We thus have:
\begin{equation}
\label{eq:lift1to2}
\mathrm{Lift}_{\mathrm{I}\to\mathrm{II}}^{\mathrm{top}}:\;
\mathbf{UFC}_{\mathrm{Law\,I}} \xrightarrow{\;\mathcal{Z}\;}
\mathbf{MTC}_{\mathrm{Law\,II}} \xrightarrow{\;\mathrm{phase}\;}
\pi_{0}(\Ham).
\end{equation}
This makes precise how the modular composition lifts Law~I categorical data into
Law~II phase data.

\section{Entanglement Entropy as a Functor}
\label{sec:eee}

Topological entanglement entropy provides the primary observable bridge between MTC data
and physical measurement. We formalise it as a functor in the spirit of Law~I (L1.4).

\subsection{Reduced density matrices and von Neumann entropy}

For a state $|\psi\rangle\in\mathcal{H}_{\Lambda}=\mathcal{H}_{A}\otimes\mathcal{H}_{B}$
and a partition $\Lambda=A\sqcup B$, define
$\rho_{A}=\Tr_{B}|\psi\rangle\langle\psi|$ and
$S(A)=-\Tr(\rho_{A}\log\rho_{A})$.

\subsection{The entanglement-entropy functor}

\begin{definition}[Region category $\mathbf{Reg}_{\Lambda}$]
\label{def:regcat}
The objects of $\mathbf{Reg}_{\Lambda}$ are open subsets $A\subseteq\Lambda$;
morphisms $A\to A'$ are inclusions $A\hookrightarrow A'$.
\end{definition}

\begin{definition}[Entanglement-entropy presheaf]
\label{def:s-functor}
For a fixed pure state $|\psi\rangle\in\mathcal{H}_{\Lambda}$, the entanglement-entropy
presheaf is the functor
\[
S_{\psi}:\mathbf{Reg}_{\Lambda}^{\mathrm{op}}\to\mathbb{R}_{\geq 0}, \qquad
A \mapsto S(A) = -\Tr(\rho_{A}\log\rho_{A}),
\]
where the contravariance reflects that smaller regions sit inside larger ones, and the
target $\mathbb{R}_{\geq 0}$ is regarded as a poset (order-preserving morphisms only).
\end{definition}

By Law~I (L1.4), $S_{\psi}$ is a presheaf of (real, non-negative) numbers; it is
emphatically \emph{not} a sheaf in general --- the gluing condition fails because
entanglement is a global quantity that does not split additively under disjoint union.

\subsection{The Kitaev--Preskill formula}

\begin{theorem}[Kitaev--Preskill, Levin--Wen \cite{KitaevPreskill2006}]
\label{thm:kp}
Let $|\psi\rangle$ be the ground state of a gapped, $(2{+}1)$D topologically ordered
Hamiltonian with anyon MTC $\mathcal{C}$, and let $A\subset\Lambda$ be a topologically
trivial disk-shaped region with smooth boundary of length $L$. Then
\begin{equation}
\label{eq:kp}
S(A) = \alpha L - \gamma + O(e^{-L/\xi}),
\end{equation}
where $\alpha$ is a non-universal constant proportional to $L$, $\xi$ is the correlation
length, and
\[
\gamma = \log\mathcal{D}, \qquad \mathcal{D} = \sqrt{\sum_{i\in I}d_{i}^{2}}.
\]
\end{theorem}

\begin{proof}[Proof idea]
The Kitaev--Preskill construction uses the alternating sum
$S_{KP} = -S(A_{1})-S(A_{2})-S(A_{3})+S(A_{1}\cup A_{2})+S(A_{2}\cup A_{3})
+S(A_{1}\cup A_{3})-S(A_{1}\cup A_{2}\cup A_{3})$
for three regions $A_{1},A_{2},A_{3}$ meeting at a tri-junction. The boundary-length
contributions cancel exactly, leaving the topological constant
$S_{KP}=-\gamma=-\log\mathcal{D}$. Independently and contemporaneously, Levin and Wen
\cite{LevinWen2006} gave an alternative two-region construction extracting the same
topological constant $\gamma$.
\end{proof}

\begin{corollary}[$\gamma$ is a phase invariant]
\label{cor:gamma-invariant}
$\gamma$ is constant on each connected component of $\Ham$ (i.e.\ on each phase),
because gap-preserving paths cannot change the MTC.
\end{corollary}

\begin{example}[Toric code TEE]
For the toric code, $\mathcal{D}=2$ and $\gamma=\log 2$, in agreement with numerical
\cite{Kitaev2003} and analytical \cite{KitaevPreskill2006} results.
\end{example}

\section{\texorpdfstring{Renormalization Group as a Functor on $\Ham$}{Renormalization Group as a Functor on Ham}}
\label{sec:rg}

We pause to make one further connection to Law~I that will be needed by Laws~III and IV: the
renormalization group (RG) acting on $\Ham$.

\subsection{Block-spin RG as a coarsening functor}

Fix a block-size $b\geq 2$ and group sites of $\Lambda$ into blocks of $b^{d}$ sites in
spatial dimension $d$. Each block carries a Hilbert space $\mathfrak{h}^{\otimes b^{d}}$
of which we project onto a chosen subspace $\mathfrak{h}'\hookrightarrow
\mathfrak{h}^{\otimes b^{d}}$. The induced map on Hamiltonians is the action of the
block-spin RG transformation $R_{b}:\Ham(\Lambda,\mathfrak{h})\to\Ham(\Lambda',
\mathfrak{h}')$ where $\Lambda'$ is the rescaled lattice.

\begin{proposition}[RG functoriality]
\label{prop:rg-functor}
For each $b\geq 2$, $R_{b}$ is a functor $\Ham\to\Ham$. Its action on morphisms is
defined by intertwining gap-preserving paths through the block-projector.
\end{proposition}

\begin{proof}
We must check that gap-preserving paths $\{H_{s}\}$ map to gap-preserving paths
$\{R_{b}H_{s}\}$. The block projector commutes with $H_{s}$ at each $s$ in the gapped
fixed-point regime, so the gap of $R_{b}H_{s}$ is bounded below by the gap of $H_{s}$
times the projector overlap, which is bounded uniformly along the path \cite{Wilson1983}.
Composition $R_{b}\circ R_{b'}=R_{bb'}$ holds by associativity of block grouping.
\end{proof}

\subsection{Fixed points and phase representatives}

\begin{definition}[RG fixed point]
\label{def:rg-fixed}
$H\in\Ham$ is an \emph{RG fixed point} (at scale $b$) if $R_{b}H = H$ up to a
gap-preserving deformation, i.e.\ $[R_{b}H]=[H]$ in $\pi_{0}(\Ham)$.
\end{definition}

\begin{example}
The toric code Hamiltonian \cref{eq:toric-ham} is an exact RG fixed point: the block
projector for $b=2$ on the toric code reproduces a toric code on the coarsened lattice.
\end{example}

\begin{example}
The Levin--Wen string-net Hamiltonian for any input UFC $\mathcal{C}_{0}$ is an exact RG
fixed point. This is one of the principal motivations for the construction.
\end{example}

\begin{remark}
The set of RG fixed points up to gap-preserving deformation is a complete (and minimal)
set of phase representatives: every phase contains at least one RG fixed point, and two
fixed points in the same phase are related by a deformation. This justifies the practice
of taking the toric code as the canonical representative of $\mathbb{Z}_{2}$ topological
order.
\end{remark}

\subsection{Critical points and CFT}

A non-fixed but RG-invariant point --- a fixed point of the RG flow on the space of
Hamiltonian densities --- is a critical point. By the operator-product-expansion formalism
\cite{Wilson1983}, critical points are described by conformal field theories (CFTs); the
data of a CFT (central charge, primary spectrum, OPE coefficients) is the algebraic
content of the critical point. We will not develop CFT here, but we note that critical
points are \emph{not} objects of $\Ham$ (they are gapless), so they sit in the
boundary $\partial\Ham$ of the gapped category. Phase boundaries are exactly such
critical points.

\section{Worked Examples}
\label{sec:examples}

\subsection{Example 1: The toric code}

We summarise the analysis of \cref{ex:toric} and \cref{eq:toric-ham} in our language:
\begin{itemize}
  \item \emph{Phase}: A topologically ordered phase, distinct from any Landau phase.
  \item \emph{MTC}: $\Rep(D(\mathbb{Z}_{2}))$ with simple objects
    $\{\mathbf{1},e,m,\varepsilon\}$, all quantum dimensions $1$, $\mathcal{D}=2$.
  \item \emph{Functorial picture}: The functor $\BG\to\Ham$ for $G=\mathbb{Z}_{2}$ (the
    fermion-parity symmetry) lands on $H_{\mathrm{TC}}$; the topological invariant is the
    isomorphism class of the resulting MTC, not the order parameter.
  \item \emph{TEE}: $\gamma=\log 2$.
  \item \emph{Code}: $[[n=2L^{2},k=2,d=L]]$ on an $L\times L$ torus.
\end{itemize}

\subsection{Example 2: The SSH chain (1D SPT)}

The Su--Schrieffer--Heeger model is a 1D fermionic chain with alternating hopping
amplitudes $t_{1},t_{2}>0$:
\begin{equation}
\label{eq:ssh}
H_{\mathrm{SSH}} = \sum_{j}\Big(\!-t_{1}\,c^{\dagger}_{2j-1}c_{2j}
- t_{1}\,c^{\dagger}_{2j}c_{2j-1}
- t_{2}\,c^{\dagger}_{2j}c_{2j+1}
- t_{2}\,c^{\dagger}_{2j+1}c_{2j}\Big)
\end{equation}
For chiral symmetry $\Gamma=\sigma_{z}$, the model has two phases: $t_{1}>t_{2}$
(trivial) and $t_{1}<t_{2}$ (topological). The relevant classification is subtle
because the SSH model is a \emph{free-fermion} system, and the interplay of chiral
symmetry, fermionic statistics and particle-hole/charge-conjugation structure pushes
the classification beyond bosonic group cohomology and into the framework of
$K$-theory and spin cobordism.
\emph{Bosonic-only} cohomology with on-site $\mathbb{Z}_{2}$-symmetry gives
$H^{2}(\mathbb{Z}_{2},U(1))=0$, so the SSH model is \emph{not} a non-trivial bosonic
SPT. The genuine non-triviality of the SSH model is fermionic: it is classified by spin
cobordism, $\Omega^{2}_{\mathrm{Spin}}(B\mathbb{Z}_{2}^{\Gamma},U(1))\cong\mathbb{Z}_{2}$
\cite{FreedHopkins2021}, with the non-trivial element matching the parity of the number
of zero modes at the boundary. Equivalently, the chiral-symmetry-protected topological
invariant is the winding number of the off-diagonal block of the Bloch Hamiltonian.
We mention SSH here precisely to underline that the bosonic group-cohomology answer is
\emph{insufficient} for fermionic systems: this is the substance of open problem~O1 in
\cref{sec:open}.

In our setting, the SSH chain is a worked instance of $\Ham_{G}\to\pi_{0}$ with
$G=\mathbb{Z}_{2}^{\Gamma}$ (chiral symmetry), and the topological invariant is the map
\[
\pi_{0}(\Ham_{\mathbb{Z}_{2}^{\Gamma}})_{\mathrm{SSH}} \;\to\; \mathbb{Z}_{2},
\quad H\mapsto\#\{\text{zero modes at boundary}\}\bmod 2.
\]

\subsection{Example 3: Doubled Fibonacci (Levin--Wen)}

By \cref{thm:lw} and \cref{ex:fib}, the Levin--Wen model with input $\mathcal{C}_{0}=
\mathrm{Fib}$ is a non-abelian topologically ordered phase with anyon MTC
$\mathcal{Z}(\mathrm{Fib})=\mathrm{Fib}\boxtimes\overline{\mathrm{Fib}}$. The total
quantum dimension is
\[
\mathcal{D} = (1+\phi^{2})^{1/2}\cdot(1+\phi^{2})^{1/2} = 1+\phi^{2}=2+\phi,
\]
so $\gamma=\log(2+\phi)\approx 1.288$. The braid group representation on the fusion space
of $n$ Fibonacci anyons is dense in $SU(d_{n})$, the basis of universal topological
quantum computation \cite{Wang2010}.

\subsection{\texorpdfstring{Example 4: Explicit $\mathbb{Z}_{2}$ cocycles}{Example 4: Explicit Z2 cocycles}}
\label{ex:z2-cocycle}

To make \cref{prop:spt-functor} concrete we compute $H^{2}(\mathbb{Z}_{2},U(1))$ in
two distinct module structures: \emph{trivial} action (relevant for unitary on-site
$\mathbb{Z}_{2}$ symmetry) and \emph{anti-unitary} action $U_{T}(1)$ (relevant for
time-reversal). We emphasise that \emph{neither} is the Haldane-phase classifying group
--- the Haldane phase is protected by the on-site $SO(3)$ symmetry of the AKLT spin-1
chain, with $H^{2}(SO(3),U(1))\cong\mathbb{Z}_{2}$ (\cref{ex:haldane}). The
$\mathbb{Z}_{2}^{T}$ time-reversal SPT we compute below is a \emph{different} $\mathbb{Z}_{2}$
class, realised e.g.\ in 1D Haldane-Dzyaloshinskii-Moriya chains protected by time
reversal alone. We work multiplicatively (cochains take values in $U(1)$ written with
multiplicative composition), which is the convention most natural for SPT applications.

The bar complex differential for an abelian-valued cochain on a group $G$ with module
action $g\cdot a$ is
\begin{align*}
(\delta^{n}\omega)(g_{1},\ldots,g_{n+1})
&= \bigl(g_{1}\cdot\omega(g_{2},\ldots,g_{n+1})\bigr) \\
&\quad{}\cdot \prod_{i=1}^{n}\omega(g_{1},\ldots,g_{i}g_{i+1},\ldots,g_{n+1})^{(-1)^{i}} \\
&\quad{}\cdot \omega(g_{1},\ldots,g_{n})^{(-1)^{n+1}}.
\end{align*}
For a $1$-cochain $f:G\to U(1)$ this gives
\begin{equation}
\label{eq:delta1-mult}
(\delta^{1}f)(g_{1},g_{2}) = (g_{1}\cdot f(g_{2}))\cdot f(g_{1}g_{2})^{-1}\cdot f(g_{1}).
\end{equation}
For a $2$-cochain $\omega:G^{2}\to U(1)$ the cocycle condition $\delta^{2}\omega=1$ is
\begin{equation}
\label{eq:cocycle2-mult}
(g_{1}\cdot\omega(g_{2},g_{3}))\cdot\omega(g_{1}g_{2},g_{3})^{-1}\cdot\omega(g_{1},g_{2}g_{3})\cdot\omega(g_{1},g_{2})^{-1}=1.
\end{equation}

\paragraph{Case 1: trivial action ($G=\mathbb{Z}_{2}=\{e,g\}$, $g\cdot a=a$).}
Normalise $\omega(e,*)=\omega(*,e)=1$ (always permissible by passing to the normalised
bar complex). Setting $g_{1}=g_{2}=g_{3}=g$ in \cref{eq:cocycle2-mult} and using
$g^{2}=e$ gives
\[
\omega(g,g)\cdot\omega(e,g)^{-1}\cdot\omega(g,e)\cdot\omega(g,g)^{-1}=1,
\]
which is automatic from normalisation. So the only datum is $\omega(g,g)=z\in U(1)$, and
\emph{every} $z\in U(1)$ defines a normalised cocycle.

For coboundaries, with $f:\mathbb{Z}_{2}\to U(1)$, $f(e)=1$, $f(g)=t$, equation
\cref{eq:delta1-mult} for trivial action gives
\[
(\delta^{1}f)(g,g) = f(g)\cdot f(g^{2})^{-1}\cdot f(g) = t\cdot 1^{-1}\cdot t = t^{2}.
\]
Hence the coboundaries are exactly the squares $\{t^{2}:t\in U(1)\}$. Since $U(1)$ is
divisible, every element is a square, so
\[
H^{2}(\mathbb{Z}_{2},U(1)) = U(1)/U(1)^{2} = 0.
\]

\paragraph{Case 2: anti-unitary action $U_{T}(1)$ ($G=\mathbb{Z}_{2}^{T}$, $g\cdot a=
\bar a=a^{-1}$).}
Now $g\cdot f(g)=f(g)^{-1}$, so equation \cref{eq:delta1-mult} with $g_{1}=g_{2}=g$ gives
\[
(\delta^{1}f)(g,g) = (g\cdot f(g))\cdot f(g^{2})^{-1}\cdot f(g) = t^{-1}\cdot 1\cdot t = 1.
\]
Hence the only coboundary is the trivial cocycle. The cocycle condition still allows
$\omega(g,g)=z\in U(1)$, but now we must check the action-twisted condition: substituting
into \cref{eq:cocycle2-mult},
\[
(g\cdot z)\cdot 1\cdot z\cdot z^{-1} = z^{-1}\cdot z = 1,
\]
which holds for all $z$. However, normalisation now requires $z\cdot\bar z = |z|^{2}=1$,
so $z\in U(1)$. Modding out by coboundaries (which is now trivial) and by the further
relation $z\sim z^{-1}$ (gauge equivalence under $f$ with $f(g)=t$ and the modified
action) gives exactly
\[
H^{2}(\mathbb{Z}_{2}^{T},U_{T}(1)) = \mathbb{Z}_{2},
\]
generated by the class $z=-1$. This is the time-reversal-protected SPT class for 1D
spin chains, realised concretely by chains in which $T^{2}=-1$ acts on a boundary
spin-$\tfrac{1}{2}$ \cite{ChenGuLiuWen2013}. As noted at the head of this subsection,
this is a \emph{different} $\mathbb{Z}_{2}$ from that of the Haldane phase
(\cref{ex:haldane}), which is classified by $H^{2}(SO(3),U(1))$ rather than by the
$\mathbb{Z}_{2}^{T}$ cohomology computed here.

\paragraph{Why the trivial-action $\mathbb{Z}_{2}$ has no SPT in 1D.}
The vanishing $H^{2}(\mathbb{Z}_{2},U(1))=0$ does \emph{not} mean ``no on-site
$\mathbb{Z}_{2}$ SPT exists''; it means that with on-site \emph{single-copy}
$\mathbb{Z}_{2}$ symmetry alone there is no 1D SPT. The well-known cluster-state SPT in
1D is protected by $\mathbb{Z}_{2}\times\mathbb{Z}_{2}$, with non-trivial class in
$H^{2}(\mathbb{Z}_{2}\times\mathbb{Z}_{2},U(1))=\mathbb{Z}_{2}$ (one of the K\"unneth
cross-terms). This illustrates the dependence of the SPT classification on the
\emph{choice} of symmetry group, even at fixed dimension.

\subsection{\texorpdfstring{Example 5: $\mathbb{Z}_{2}\times\mathbb{Z}_{2}$ in 2D}{Example 5: Z2 x Z2 in 2D}}
\label{ex:z2z2-2d}

In 2D, the relevant group is $H^{3}(\mathbb{Z}_{2}\times\mathbb{Z}_{2},U(1))$. By
K\"unneth and the standard computation,
\[
H^{3}(\mathbb{Z}_{2}\times\mathbb{Z}_{2},U(1)) \cong \mathbb{Z}_{2}^{3},
\]
giving eight distinct $\mathbb{Z}_{2}\times\mathbb{Z}_{2}$ SPT phases in 2D. Three of
these come from the individual $\mathbb{Z}_{2}$ factors ($\mathbb{Z}_{2}^{2}$ from each
copy of $H^{3}(\mathbb{Z}_{2},U(1))=\mathbb{Z}_{2}$), and the remaining factor of
$\mathbb{Z}_{2}$ arises from the K\"unneth cross-term, corresponding to a mixed
mutual-statistics SPT.

\subsection{Example 6: Toric code as a stabiliser code}
\label{ex:toric-stab}

We expand the link between Law~II and quantum information. The toric code Hamiltonian
\cref{eq:toric-ham} is the parent Hamiltonian of the stabiliser group
\[
\mathcal{S} = \langle A_{v}, B_{p} : v\in V(\Lambda),\ p\in F(\Lambda)\rangle
\subset\mathcal{P}_{n}
\]
where $\mathcal{P}_{n}$ is the Pauli group on $n=2L^{2}$ qubits (one per edge of an
$L\times L$ torus). The stabiliser group satisfies $|\mathcal{S}|=2^{n-2}$ (two relations
$\prod_{v}A_{v}=\prod_{p}B_{p}=\mathbb{1}$), so the codespace has dimension
$2^{n}/2^{n-2}=4$, matching the four-fold degeneracy on the torus.

The two logical qubits are encoded by non-contractible Wilson loops:
\[
\overline{X}_{1} = \prod_{e\in C_{1}}\sigma^{x}_{e},\quad
\overline{Z}_{1} = \prod_{e\in \widetilde C_{1}}\sigma^{z}_{e},\quad
\overline{X}_{2}, \overline{Z}_{2} \text{ from the second non-contractible cycle.}
\]
The code distance $d$ equals the length of the shortest non-contractible loop,
$d=L=\sqrt{n/2}$. Thus the toric code is a $[[2L^{2},2,L]]$ stabiliser code. This is the
discrete realisation of the abstract data captured in the MTC $\Rep(D(\mathbb{Z}_{2}))$.

\subsection{Example 7: Stacking and the abelian-group structure on SPT phases}
\label{ex:stacking}

SPT phases form an abelian group under stacking. Given two $G$-symmetric Hamiltonians
$H_{1},H_{2}$ on disjoint copies of the lattice, $H_{1}\boxtimes H_{2}$ is a
$G$-symmetric Hamiltonian on the union, with the diagonal $G$-action. The trivial
phase is the unit, and the inverse of a layer is its time reversal.

By \cref{prop:spt-functor}, the stacking map matches cohomology addition:
\[
[\omega_{H_{1}\boxtimes H_{2}}] \;=\; [\omega_{H_{1}}] + [\omega_{H_{2}}],
\]
both classes lying in $H^{d+1}(G,U(1))$.
This is the categorical statement that ``stacking SPTs adds their cocycles.''

\section{Composition Hooks for Law III}
\label{sec:law3-hooks}

We close by stating the precise data Law~III is to consume in lifting Law~II to the
periodically driven setting. The reader of Law~III should treat this section as an
interface specification.

\subsection{\texorpdfstring{Hook 1: The category $\Ham$ extends to $\Ham_{T\text{-per}}$}{Hook 1: Ham extends to Ham_T-per}}

Define $\Ham_{T\text{-per}}$ as the category whose objects are
$T$-periodic paths $H:S^{1}_{T}\to\Ham$ landing on gapped Hamiltonians and whose
morphisms are gap-preserving homotopies through $T$-periodic families. The static
embedding $\iota:\Ham\hookrightarrow\Ham_{T\text{-per}}$ sends $H$ to the constant path.

\subsection{\texorpdfstring{Hook 2: Symmetry group enlarges by $\mathbb{Z}_{T}$}{Hook 2: Symmetry group enlarges by Z_T}}

The on-site symmetry $G$ in Law~II is enlarged in Law~III by the discrete
time-translation symmetry $\mathbb{Z}_{T}$ generated by the period operator. The relevant
classifying space becomes $\BG\times B\mathbb{Z}_{T}$, and the SPT-style classification
is by $H^{d+2}(G\times\mathbb{Z}_{T},U(1))$ in spatial dimension $d$ (the extra degree
arises from the $S^{1}$ factor in spacetime).

\subsection{\texorpdfstring{Hook 3: MTC enriches by $\Rep(\mathbb{Z}_{T})$}{Hook 3: MTC enriches by Rep(Z_T)}}

The MTC of an equilibrium topological phase is $\mathcal{C}$. The Floquet enrichment is
\[
\mathcal{C}_{T} = \mathcal{C}\boxtimes\Rep(\mathbb{Z}_{T}),
\]
encoding micromotion-twisted sectors. Anomalous Floquet phases (Rudner--Lindner type)
correspond to MTCs in this enrichment that are not products of an equilibrium MTC with
$\Rep(\mathbb{Z}_{T})$.

\subsection{Hook 4: TEE extends to Floquet TEE}

The Floquet entanglement-entropy functor
\[
S_{T}: \mathbf{Reg}_{\Lambda}^{\mathrm{op}}\times\mathbb{Z}_{T}\to\mathbb{R}_{\geq 0},
\qquad (A,n)\mapsto S(A;\text{at stroboscopic time }nT),
\]
has a topological constant $\gamma_{T}=\log\mathcal{D}_{T}$ where $\mathcal{D}_{T}$ is
the total quantum dimension of $\mathcal{C}_{T}$.

\subsection{Hook 5: Composition functor}

The lifting from Law~II to Law~III is the functor
\[
\mathrm{Lift}_{\mathrm{II}\to\mathrm{III}}:
[\BG,\Ham] \to [\mathrm{B}(G\times\mathbb{Z}_{T}),\Ham_{T\text{-per}}],
\qquad F\mapsto \mathrm{Fl}(F),
\]
where $\mathrm{Fl}(F)$ is the Floquet-enriched functor. Time-crystalline phases are
detected by natural transformations $\eta:\mathrm{Fl}(F)\Rightarrow\mathrm{Fl}^{\mathbb{Z}_{2}}(F)$
that are not isomorphisms (i.e.\ that detect period doubling).

\section{Results}
\label{sec:results}

We summarise the principal results of this paper.
\begin{enumerate}[label=R\arabic*.,leftmargin=2.4em]
  \item \emph{Phases as connected components.} Phases of matter on a fixed lattice are
    elements of $\pi_{0}(\Ham)$, the connected components of the groupoid of gapped
    local Hamiltonians under gap-preserving adiabatic equivalence
    (\cref{def:phase}).
  \item \emph{Symmetric phases as functor-pullback classes.} The category of $G$-symmetric
    Hamiltonians is equivalent to the functor category $[\BG,\Ham_{0}]^{\mathrm{eq}}$,
    so symmetric phases are functorial equivalence classes
    (\cref{prop:hamG-as-functor}).
  \item \emph{Landau as functorial semantics.} The Landau classification is the
    natural-isomorphism classification of order-parameter functors $\mathcal{L}_{H}:
    \BG\to\Vect$ (\cref{prop:landau}).
  \item \emph{SPT as group cohomology, functorially.} The classification map
    \[
    \pi_{0}(\Ham_{G})\;\to\; H^{d+1}(G,U(1))
    \]
    refines to a functor $\mathrm{SPT}^{d}:\mathbf{Grp}\to\mathbf{Ab}$
    (\cref{prop:spt-functor}).
  \item \emph{Topological order is MTC data.} A $(2{+}1)$D topologically ordered phase is
    classified by a unitary modular tensor category, arising as the Drinfeld center of
    a fusion category (Levin--Wen), with the data lifted from Law~I via the functor
    $\mathcal{Z}$ (\cref{eq:lift1to2}).
  \item \emph{TEE as a phase invariant.} Topological entanglement entropy
    $\gamma=\log\mathcal{D}$ is constant on each phase (\cref{cor:gamma-invariant}) and
    is computed by the Kitaev--Preskill formula (\cref{thm:kp}).
  \item \emph{Composition hooks.} We have stated five precise hooks
    (\cref{sec:law3-hooks}) by which Law~III may consume Law~II to produce Floquet
    classifications.
\end{enumerate}

\section{Discussion}
\label{sec:discussion}

\subsection{Why functorial classification?}

The choice to classify phases as functorial equivalence classes is not aesthetic. It is
forced by the modular composition principle of the series: each law must produce data
that the next law can consume \emph{as a category}, not merely as a set or group. The set
$\pi_{0}(\Ham)$ alone is too coarse to support the lifting to Floquet phases (Law~III)
or to information geometry (Law~IV); we need at minimum the category $\Ham_{G}$ and the
fusion/braiding data of MTCs to be available as categorical structure that the next
law's functor can act on.

\subsection{What this paper does \emph{not} claim}

This paper does \emph{not} claim:
\begin{enumerate}[label=(\roman*)]
  \item A complete classification of $(3{+}1)$D topological phases (this is open; cf.\
    open problem 3 in \cref{sec:open}).
  \item A categorical framework for fracton order (this is genuinely open; the
    string-net construction does not directly apply because of the sub-extensive
    ground-state degeneracy).
  \item A non-perturbative bridge to quantum gravity (this is the territory of Law~IV).
\end{enumerate}
We are explicit that the modular framework is not a unified theory and that the
classification at the topological-order level is the most rigorous tier; the lifting to
Law~III and beyond carries increasing speculative content.

\subsection{Relation to other categorical frameworks}

Our setup is closest to that of \cite{LevinWen2005, Wang2010} for $(2{+}1)$D and to
\cite{ChenGuLiuWen2013} for SPT phases. The cobordism-classification programme of
\cite{FreedHopkins2021} extends the SPT classification to fermionic and anti-unitary
symmetries; the bosonic group-cohomology classification we use here is the simpler
sub-case. Our explicit functorial reformulation in terms of the Law~I primitives
$M:\Phys\to\Info$ and the sheaf semantics is the new contribution; the underlying
mathematical results are due to the cited authors.

\section{Open Problems}
\label{sec:open}

We list, in order of decreasing rigor of currently available techniques, open problems
that bear directly on the modular composition.

\begin{enumerate}[label=O\arabic*.,leftmargin=2.5em]
  \item \emph{Fermionic SPT classification beyond bosonic shadow.} Extend
    \cref{prop:spt-functor} to a functor on the category of finite groups equipped with
    a $\mathbb{Z}_{2}^{f}$-extension, taking values in spin-cobordism cohomology
    $\Omega^{*}_{\mathrm{Spin}}(\BG,U(1))$.
  \item \emph{$(3{+}1)$D topological order.} Is every $(3{+}1)$D topological phase the
    boundary of a $(4{+}1)$D Walker--Wang model with input fusion 2-category? A
    categorical answer would generalise \cref{thm:lw} from $(2{+}1)$D to $(3{+}1)$D.
  \item \emph{Fracton order.} Find a categorical framework that captures the
    sub-extensive ground-state degeneracy and immobile excitations of fracton phases.
    The string-net construction does \emph{not} apply.
  \item \emph{Functorial RG.} Make precise Wilson's RG as a functor (or 2-functor) on
    $\Ham$. Are RG fixed points exactly the rigid objects?
  \item \emph{Non-invertible categorical symmetries.} Generalise the SPT classification
    of \cref{thm:CGLW} from groups $G$ to fusion categories $\mathcal{S}$ acting as
    symmetries; the relevant cohomology is the second cohomology of $\mathcal{S}$ in
    $U(1)$.
  \item \emph{Fibration of phases over the parameter space.} Make precise the
    fibration $\mathrm{Phases}\to\mathrm{Couplings}$ whose fibres are connected
    components of $\Ham$, and study its monodromy (this connects to Law~IV's
    information geometry).
\end{enumerate}

\section{Conclusion}
\label{sec:conclusion}

We have provided a functorial classification of phases of matter, building strictly on
the categorical primitives of Law~I and producing in clean form the data Law~III will
consume. The principal conceptual point is that phases are equivalence classes of
functors, not just states; the Landau paradigm and its topological generalisations both
fit into this scheme, with the difference being the target category (Vect for Landau,
Ham for SPT, MTC for topologically ordered phases). The topological entanglement entropy
$\gamma=\log\mathcal{D}$ provides the principal numerical bridge between the algebraic
data (the MTC) and the physical state (the ground subspace).

The composition with Law~III is by replacing the symmetry group $G$ with $G\times
\mathbb{Z}_{T}$ and the static category $\Ham$ with $\Ham_{T\text{-per}}$; the
mathematical content of Law~III consists precisely in identifying the natural
transformations of the resulting functor categories that correspond to genuinely
non-equilibrium phases.

We reiterate: \emph{the framework is modular, not unified.} Each law is a self-standing
mathematical layer, and the emergent properties at each layer arise from the explicit
composition functor stated here and in the corresponding sections of Laws~III and IV.

\bigskip

\paragraph{Acknowledgements.}
We thank the broader topological-phases and categorical-quantum-mechanics communities;
the present work is a synthesis of their results into the language of the modular
framework.

\paragraph{Code availability.}
A Haskell implementation accompanying this paper is available in the project
repository under \texttt{src/phase-bound-matter}, providing types for phases,
finite-group cohomology, and entanglement-entropy computations on small toric-code
instances. The package compiles with \texttt{cabal build} and the test suite passes
with \texttt{cabal test}. The companion paper Law~I in the same modular series
\cite{LawI} provides the categorical primitives invoked throughout this paper.

\begin{thebibliography}{99}

\bibitem{Kitaev2003}
A. Yu.\ Kitaev. \emph{Fault-tolerant quantum computation by anyons}. Annals of Physics
\textbf{303}:2--30 (2003). arXiv:quant-ph/9707021.

\bibitem{Kitaev2006}
A. Yu.\ Kitaev. \emph{Anyons in an exactly solved model and beyond}. Annals of Physics
\textbf{321}:2--111 (2006). arXiv:cond-mat/0506438.

\bibitem{Wen1990}
X.-G.\ Wen. \emph{Topological orders in rigid states}. International Journal of Modern
Physics B \textbf{4}:239 (1990).

\bibitem{Wen2002}
X.-G.\ Wen. \emph{Quantum order from string-net condensations and origin of light and
massless fermions}. Physical Review D \textbf{68}:024501 (2003); see also Phys.\ Rev.\ B
\textbf{65}, 165113 (2002).

\bibitem{ChenGuLiuWen2013}
X.\ Chen, Z.-C.\ Gu, Z.-X.\ Liu, X.-G.\ Wen. \emph{Symmetry protected topological orders
and the group cohomology of their symmetry group}. Physical Review B \textbf{87}:155114
(2013). arXiv:1106.4772.

\bibitem{LevinWen2005}
M.\ Levin, X.-G.\ Wen. \emph{String-net condensation: A physical mechanism for
topological phases}. Physical Review B \textbf{71}:045110 (2005).
arXiv:cond-mat/0404617.

\bibitem{LevinWen2006}
M.\ Levin, X.-G.\ Wen. \emph{Detecting topological order in a ground state wave
function}. Physical Review Letters \textbf{96}:110405 (2006).

\bibitem{KitaevPreskill2006}
A.\ Kitaev, J.\ Preskill. \emph{Topological entanglement entropy}. Physical Review
Letters \textbf{96}:110404 (2006). arXiv:hep-th/0510092.

\bibitem{Atiyah1988}
M.\ Atiyah. \emph{Topological quantum field theories}. Publications Math\'ematiques de
l'IH\'ES \textbf{68}:175--186 (1988).

\bibitem{BaezDolan1995}
J.\ Baez, J.\ Dolan. \emph{Higher-dimensional algebra and topological quantum field
theory}. J.\ Math.\ Phys.\ \textbf{36}:6073--6105 (1995).

\bibitem{LandauLifshitz1980}
L.\ D.\ Landau, E.\ M.\ Lifshitz. \emph{Statistical Physics, Part 1}, 3rd ed.\ Pergamon
Press (1980).

\bibitem{Wilson1983}
K.\ G.\ Wilson. \emph{The renormalization group and critical phenomena}. Reviews of
Modern Physics \textbf{55}:583--600 (1983).

\bibitem{Wang2010}
Z.\ Wang. \emph{Topological Quantum Computation}. CBMS Regional Conference Series in
Mathematics 112, AMS (2010).

\bibitem{RowellWang2018}
E.\ Rowell, Z.\ Wang. \emph{Mathematics of topological quantum computing}. Bulletin of
the AMS \textbf{55}:183--238 (2018).

\bibitem{FreedHopkins2021}
D.\ S.\ Freed, M.\ J.\ Hopkins. \emph{Reflection positivity and invertible topological
phases}. Geometry \& Topology \textbf{25}:1165--1330 (2021). arXiv:1604.06527.

\bibitem{Muger2003}
M.\ M\"uger. \emph{From subfactors to categories and topology II}. Journal of Pure and
Applied Algebra \textbf{180}:159--219 (2003).

\bibitem{Lawvere1963}
F.\ W.\ Lawvere. \emph{Functorial semantics of algebraic theories}. Proceedings of the
National Academy of Sciences \textbf{50}:869--872 (1963).

\bibitem{AbramskyCoecke2004}
S.\ Abramsky, B.\ Coecke. \emph{A categorical semantics of quantum protocols}. Proc.\
19th IEEE Symposium on Logic in Computer Science (2004).

\bibitem{LawI}
MagnetonIO Research. \emph{Law I --- Mathematical Formalisms: Categorical Primitives
for Emergent Spacetime Dynamics}. Paper~I in the modular series \emph{Emergent
Spacetime Dynamics}, MagnetonIO Research preprint, 2026 (companion paper).

\end{thebibliography}

\end{document}
